% Complete Experimental Table - Docker vs Apple container Benchmark
% Automatically generated by analyze_results_enhanced.py
%
% To use in LaTeX, copy the contents of the desired section
%
\documentclass{article}
\usepackage{booktabs}
\usepackage{longtable}
\usepackage{multirow}
\usepackage{array}
\begin{document}

% ============================================
% TABLE: CPU/Memory
% ============================================

\begin{table}[h]
\centering
\small
\caption{Experimental Results: CPU/Memory}
\label{tab:cpu_memory}
\begin{tabular}{lccccccccc}
\toprule
Scenario & Tool & N & Mean & SD & Median & IQR & CI 95\% & Effect Size & p-value (FDR) \\
 & & & & & & & [lower, upper] & (Cohen's d) & \\
\midrule
\multirow{2}{*}{cpu\_1t} & Docker & 30 & 8771.6943 & 17.1922 & 8773.3350 & 18.0650 & [8765.2747, 8778.1140] & -1.0257 & 0.0002*** \\
 & Container & 30 & 8789.1940 & 16.9307 & 8790.9550 & 16.3675 & [8782.8720, 8795.5160] &  &  \\
\midrule
\multirow{2}{*}{cpu\_2t} & Docker & 30 & 16819.9450 & 42.3729 & 16814.4750 & 44.8700 & [16804.1227, 16835.7673] & -1.5697 & 0.0000*** \\
 & Container & 30 & 16881.1223 & 35.2474 & 16876.8050 & 33.9125 & [16867.9607, 16894.2839] &  &  \\
\midrule
\bottomrule
\end{tabular}
\begin{tablenotes}
\small
\item N: sample size; SD: standard deviation; IQR: interquartile range; CI: confidence interval; Effect Size: effect size (Cohen's d); p-value corrected by FDR (Benjamini-Hochberg).
\item Significance: *** p < 0.001, ** p < 0.01, * p < 0.05
\end{tablenotes}
\end{table}

\end{document}
